% TOP_PROBLEM: {"id": 30005169, "title": "Known Cases of Iitaka Subadditivity", "category_id": 5, "category": {"id": 5, "name": "algebraic_geometry", "display_name": "Algebraic Geometry", "description": "Geometric objects defined by polynomial equations.", "slug": "algebraic-geometry", "order_index": 5, "created_at": "2026-07-31T15:26:25.671Z"}, "status": "partially_solved"}
% FIRST_POSTED: 2026-10-06
% ATTEMPT_STATUS: unresolved
% OPENAI_CLAIM: {"schema_version":1,"source_repo":"https://github.com/openai/math","source_commit":"adc7f1241b42e322a6451854ab7e4b4c146bf78a","release_date":"2026-10-06","match":"partial","resolution":"partial","independently_reviewed":false,"second_pass":"n/a","families":[{"family":"033","title":"Iitaka subadditivity, variation, and logarithmic additivity","subject":"Algebraic and complex geometry","manuscripts":[{"title":"Logarithmic Kodaira dimension and whole-fiber variation","pdf_path":"preprints/Logarithmic-Kodaira-dimension-and-whole-fiber-variation-September-26-2026/paper.pdf","readme_path":"preprints/Logarithmic-Kodaira-dimension-and-whole-fiber-variation-September-26-2026/README.md","date":"2026-09-26","theorem_ref":"build/sections/01-introduction.tex:38, Theorem thm:main"},{"title":"Orbifold and logarithmic Iitaka subadditivity","pdf_path":"preprints/Orbifold-and-logarithmic-Iitaka-subadditivity-September-26-2026/paper.pdf","readme_path":"preprints/Orbifold-and-logarithmic-Iitaka-subadditivity-September-26-2026/README.md","date":"2026-09-26","theorem_ref":"build/sections/introduction.tex:135, Theorem main:orbifold"},{"title":"Projective Hodge lines and ordinary Iitaka subadditivity","pdf_path":"preprints/Projective-Hodge-lines-and-ordinary-Iitaka-subadditivity-September-27-2026/paper.pdf","readme_path":"preprints/Projective-Hodge-lines-and-ordinary-Iitaka-subadditivity-September-27-2026/README.md","date":"2026-09-27","theorem_ref":"build/sections/03-limits.tex:19, Theorem altord:h:inputs"}],"lean":{"doc_path":"lean/docs/033.md","comparators":[],"declarations":[],"covers_main_theorem":false},"reasoning_trace":null}]}
\documentclass[11pt]{article}
\usepackage[margin=25mm]{geometry}
\usepackage{fontspec}
\setmainfont{Latin Modern Roman}
\usepackage{amsmath,amssymb,mathtools}
\usepackage{unicode-math}
\setmathfont{Latin Modern Math}
\usepackage{enumitem}
\usepackage{xurl,hyperref,bookmark}
\hypersetup{colorlinks=true,urlcolor=blue}
\setcounter{secnumdepth}{0}
\setlength{\parindent}{0pt}
\setlength{\parskip}{5pt}
\setlength{\emergencystretch}{3em}
\begin{document}
\section*{Known Cases of Iitaka Subadditivity}

\subsection{OpenAI's claim}
Manuscript 1 (family 033, 2026-09-26): The logarithmic Iitaka--Viehweg lower bound is claimed for projective fibrations over a smooth quasi-projective base of nonnegative logarithmic Kodaira dimension, using variation of the entire generic fiber. Manuscript 2 (family 033, 2026-09-26): For compact manifolds in Fujiki class $\mathcal C$ carrying rational normal-crossing boundaries, the orbifold Iitaka lower bound is claimed, implying its ordinary and logarithmic specializations. Manuscript 3 (family 033, 2026-09-27): The ordinary Kodaira dimension inequality is claimed for projective connected-fiber morphisms in characteristic zero, without restricting total dimension or assuming nonnegative Kodaira dimension upstairs.

\subsection{Outline of the argument}
Manuscript 1 (family 033, 2026-09-26): Canonical root covers isolate discriminant and moduli terms without losing pluriform orders. Flatness of the relevant Hodge lines is converted into birational constancy by singular-weight extension estimates and a big adjoint model; marked-family positivity and finite-cover descent then account for whole-fiber variation.

Manuscript 2 (family 033, 2026-09-26): A relative Iitaka reduction first replaces the original fibers by the part for which general-type addition applies. Root-cover comparisons retain the boundary multiplicities, and positivity of the resulting period line supplies the missing contribution from the base.

Manuscript 3 (family 033, 2026-09-27): A canonical root cover of a Kodaira-zero fiber gives a one-dimensional top Hodge space. The argument distinguishes the full period image from the rank seen by that line, establishes the required positivity on the compactification, and applies the canonical bundle formula to recover subadditivity.

\subsection{Scope relative to this problem}
Manuscript 1 (family 033, 2026-09-26): The logarithmic Iitaka--Viehweg lower bound is claimed for projective fibrations over a smooth quasi-projective base of nonnegative logarithmic Kodaira dimension, using variation of the entire generic fiber. Complete coverage of the precise assertion here is not established in this catalogue; this entry retains partial support.

Manuscript 2 (family 033, 2026-09-26): For compact manifolds in Fujiki class $\mathcal C$ carrying rational normal-crossing boundaries, the orbifold Iitaka lower bound is claimed, implying its ordinary and logarithmic specializations. Complete coverage of the precise assertion here is not established in this catalogue; this entry retains partial support.

Manuscript 3 (family 033, 2026-09-27): The ordinary Kodaira dimension inequality is claimed for projective connected-fiber morphisms in characteristic zero, without restricting total dimension or assuming nonnegative Kodaira dimension upstairs. Complete coverage of the precise assertion here is not established in this catalogue; this entry retains partial support.

\subsection{Formal verification}
For family 033, the scope page documents the following selections. Only the case of a negative-infinite logarithmic fiber dimension is selected; it forces negative-infinite dimension upstairs and vanishing of positive pluriforms. Finite-dimensional fibers and the negative-base case remain outside the linked statement. Manuscript 1 is not among the papers selected by that page, so the family link does not formalize its main conclusion. Manuscript 2 is not among the papers selected by that page, so the family link does not formalize its main conclusion. Manuscript 3 is not among the papers selected by that page, so the family link does not formalize its main conclusion.

This site has not executed the Lean files or checked their dependencies. These descriptions report the scope of the linked declarations, rather than a new verification of the proofs.

\subsection{Status caveat}
OpenAI attributes these results to an internal model. Its release notice cautions that results lacking full formalization can contain mistakes. This site has not independently assessed the proofs; a solved label records the scope of the reported claim.

\subsection{Sources}
\begin{itemize}
\item OpenAI, \emph{Logarithmic Kodaira dimension and whole-fiber variation}, 2026-09-26; citation key \texttt{OAI:Logarithmic-Kodaira-dimension-and-whole-fiber-variation-September-26-2026}. \href{https://github.com/openai/math/blob/adc7f1241b42e322a6451854ab7e4b4c146bf78a/preprints/Logarithmic-Kodaira-dimension-and-whole-fiber-variation-September-26-2026/paper.pdf}{PDF (GitHub)}; \href{https://github.com/openai/math/raw/adc7f1241b42e322a6451854ab7e4b4c146bf78a/preprints/Logarithmic-Kodaira-dimension-and-whole-fiber-variation-September-26-2026/paper.pdf}{PDF (direct)}; \href{https://github.com/openai/math/blob/adc7f1241b42e322a6451854ab7e4b4c146bf78a/preprints/Logarithmic-Kodaira-dimension-and-whole-fiber-variation-September-26-2026/README.md}{Manuscript page and citation}; \href{https://github.com/openai/math/blob/main/preprints/Logarithmic-Kodaira-dimension-and-whole-fiber-variation-September-26-2026/paper.pdf}{Latest version}.
\item OpenAI, \emph{Orbifold and logarithmic Iitaka subadditivity}, 2026-09-26; citation key \texttt{OAI:Orbifold-and-logarithmic-Iitaka-subadditivity-September-26-2026}. \href{https://github.com/openai/math/blob/adc7f1241b42e322a6451854ab7e4b4c146bf78a/preprints/Orbifold-and-logarithmic-Iitaka-subadditivity-September-26-2026/paper.pdf}{PDF (GitHub)}; \href{https://github.com/openai/math/raw/adc7f1241b42e322a6451854ab7e4b4c146bf78a/preprints/Orbifold-and-logarithmic-Iitaka-subadditivity-September-26-2026/paper.pdf}{PDF (direct)}; \href{https://github.com/openai/math/blob/adc7f1241b42e322a6451854ab7e4b4c146bf78a/preprints/Orbifold-and-logarithmic-Iitaka-subadditivity-September-26-2026/README.md}{Manuscript page and citation}; \href{https://github.com/openai/math/blob/main/preprints/Orbifold-and-logarithmic-Iitaka-subadditivity-September-26-2026/paper.pdf}{Latest version}.
\item OpenAI, \emph{Projective Hodge lines and ordinary Iitaka subadditivity}, 2026-09-27; citation key \texttt{OAI:Projective-Hodge-lines-and-ordinary-Iitaka-subadditivity-September-27-2026}. \href{https://github.com/openai/math/blob/adc7f1241b42e322a6451854ab7e4b4c146bf78a/preprints/Projective-Hodge-lines-and-ordinary-Iitaka-subadditivity-September-27-2026/paper.pdf}{PDF (GitHub)}; \href{https://github.com/openai/math/raw/adc7f1241b42e322a6451854ab7e4b4c146bf78a/preprints/Projective-Hodge-lines-and-ordinary-Iitaka-subadditivity-September-27-2026/paper.pdf}{PDF (direct)}; \href{https://github.com/openai/math/blob/adc7f1241b42e322a6451854ab7e4b4c146bf78a/preprints/Projective-Hodge-lines-and-ordinary-Iitaka-subadditivity-September-27-2026/README.md}{Manuscript page and citation}; \href{https://github.com/openai/math/blob/main/preprints/Projective-Hodge-lines-and-ordinary-Iitaka-subadditivity-September-27-2026/paper.pdf}{Latest version}.
\item Family 033: \href{https://github.com/openai/math/blob/adc7f1241b42e322a6451854ab7e4b4c146bf78a/lean/docs/033.md}{Lean scope}.
\item \href{https://github.com/openai/math/blob/adc7f1241b42e322a6451854ab7e4b4c146bf78a/CONTENTS.md}{OpenAI catalogue entry}.
\end{itemize}
\end{document}
